\chapter{Literature and Technical Background}
\label{ch:literature}

\section{Photovoltaic Energy in Morocco}

Morocco receives between 2,800 and 3,400 hours of sunshine per year, with average global horizontal irradiance (GHI) ranging from 4.5 to 5.5~kWh/m$^2$/day depending on the region \cite{morocco_energy}. The city of Meknès, located in north-central Morocco at latitude 33.9°N, receives approximately 5.0~kWh/m$^2$/day annual average GHI, making it suitable for rooftop PV installations.

The Moroccan government has implemented several programs to promote distributed solar generation, including net metering regulations for installations up to 5~MW and incentives for self-consumption.

\section{Battery Energy Storage Systems}

Lithium iron phosphate (LiFePO4) batteries have become the standard for residential energy storage due to their high cycle life ($>$4,000 cycles at 80\% depth of discharge), thermal stability, and declining costs \cite{lifep04}. The state of charge (SOC) of a battery is governed by:

\begin{equation}
    \text{SOC}(t+1) = \text{SOC}(t) + \frac{\eta_{\text{ch}} \cdot P_{\text{ch}}(t) \cdot \Delta t}{E_{\text{bat}}} - \frac{P_{\text{dis}}(t) \cdot \Delta t}{\eta_{\text{dis}} \cdot E_{\text{bat}}}
    \label{eq:soc}
\end{equation}

where $\eta_{\text{ch}}$ and $\eta_{\text{dis}}$ are the charging and discharging efficiencies, $P_{\text{ch}}$ and $P_{\text{dis}}$ are the charging and discharging powers, $\Delta t$ is the time step, and $E_{\text{bat}}$ is the nominal battery capacity.

\section{Energy Management Systems}

Energy management systems for buildings with distributed generation can be classified into three categories \cite{ems_review}:

\begin{enumerate}
    \item \textbf{Rule-based systems:} Use predefined if-then rules to manage energy flows. Simple to implement but cannot anticipate future conditions.

    \item \textbf{Optimization-based systems:} Formulate energy dispatch as a mathematical optimization problem (LP, MILP, or nonlinear programming) and solve it over a finite horizon using forecasted data.

    \item \textbf{Learning-based systems:} Use reinforcement learning or other adaptive algorithms to learn dispatch policies from experience. Require extensive training data and offer less interpretability.
\end{enumerate}

This project implements and compares approaches (1) and (2).

\section{Machine Learning for Energy Forecasting}

Short-term load forecasting (STLF) for buildings typically uses features derived from historical consumption, calendar variables (hour, day of week, season), and weather data (temperature, irradiance). Common models include:

\begin{itemize}[nosep]
    \item \textbf{Linear Regression:} Simple, interpretable baseline.
    \item \textbf{Random Forest:} Ensemble of decision trees; handles nonlinearity without extensive tuning.
    \item \textbf{XGBoost:} Gradient-boosted trees; state-of-the-art for tabular data \cite{xgboost}.
    \item \textbf{LSTM networks:} Recurrent neural networks for sequence modeling; justified only when large datasets are available.
\end{itemize}

Studies on residential STLF report typical MAPE values between 5\% and 20\%, depending on building type, data quantity, and forecast horizon \cite{sklearn}.

\section{Anomaly Detection in Energy Systems}

Anomaly detection methods for energy consumption include:

\begin{itemize}[nosep]
    \item \textbf{Statistical thresholds:} Flag values exceeding $\mu \pm k\sigma$ for each time-of-day bin. Simple and interpretable.
    \item \textbf{Isolation Forest:} Unsupervised algorithm that isolates anomalies by recursive random partitioning \cite{isolationforest}. Effective for multivariate data.
\end{itemize}

\section{Mixed Integer Linear Programming for Energy Dispatch}

The energy dispatch problem is naturally formulated as a MILP when binary decisions are involved (e.g., battery charge/discharge mode, flexible load on/off). The PuLP library provides a Python interface to the CBC open-source solver \cite{pulp}, which handles problems with hundreds of variables and constraints in under one second.

\section{IoT Architecture for Energy Monitoring}

MQTT (Message Queuing Telemetry Transport) is the de facto standard for IoT communication in energy monitoring applications \cite{mqtt}. It provides a lightweight publish/subscribe protocol suitable for constrained devices such as the ESP32 microcontroller \cite{esp32}.
